\documentclass[12pt,a4paper]{article}

% Pacotes essenciais
\usepackage[utf8]{inputenc}
\usepackage[T1]{fontenc}
\usepackage[portuguese]{babel}
\usepackage{amsmath,amsfonts,amssymb}
\usepackage{graphicx}
\usepackage[a4paper,margin=2cm]{geometry}

% Título e cabeçalho
\title{{\footnotesize Universidade Federal de Minas Gerais\\Escola de Engenharia\\ELE077 Otimização Não Linear\\} \vspace{5mm} Trabalho Computacional I}

% Autores do trabalho
\author{Otávio Serafim de Souza Matos \and Nome Completo 2}
\date{\today}

\begin{document}

\maketitle

\section{Problema 1}

Resolução do primeiro problema apresentado no trabalho.
Ao analisar a função: $$\omega^{2} = R(x) = \frac{x^{T}Kx}{x^{T}Mx}$$ é possível observar que ela não é estritamente quadrática mas sim uma razão entre duas formas quadráticas. Por esse fator, e por ser uma função contínua com poucas variáveis demos prioridade ao método de Quasi-Newton para resolução do problema.

A função objetivo foi implementada utilizando operações matriciais junto da biblioteca scipy para realizar a minimização. Executando o código, obtivemos os seguintes resultados:
\begin{figure}[h]
    \centering
    \includegraphics[width=0.85\textwidth]{Figure_1.png}
    \caption{Contorno do quociente de Rayleigh com trajetória do BFGS (esquerda) e convergência de R(x) (direita).}
    \label{fig:problema1}
\end{figure}
\begin{verbatim}
Vetor x ótimo: [0.60837422 1.09624041 1.36697757]
Função-objetivo final R(x*): 0.19806226425143267
Número de iterações: 4
Número de avaliações: 20
\end{verbatim}
O valor final da função-objetivo, $R(\mathbf{x}^*) \approx 0.1981$, corresponde ao autovalor mínimo do sistema. Fisicamente, este valor representa o quadrado da frequência natural fundamental ($\omega_1^2$). Portanto, a primeira frequência natural do sistema é $\omega_1 \approx 0.445$ rad/s.
Do ponto de vista computacional, o método de Quasi-Newton apresentou alta eficiência, convergindo para o ponto ótimo em apenas 4 iterações e 20 avaliações da função-objetivo. O gráfico de convergência corrobora esse desempenho, demonstrando uma queda acentuada do valor de $R(\mathbf{x})$ já nas iterações iniciais, seguida de rápida estabilização no mínimo global esperado para a configuração física do problema. A escolha do método se provou perfeitamente adequada para a topologia da função.


\section{Problema 2}

Resolução do segundo problema apresentado no trabalho.

\section{Problema 3}

Resolução do terceiro problema apresentado no trabalho.

\section{Problema 4}

Resolução do quarto problema apresentado no trabalho.

\end{document}